\chapter{МЕТОД ОГРАНИЧЕННОЙ ГЕНЕРАЦИИ МОДЕЛЕЙ}

\section{Общая схема конвейера}

Предлагаемый конвейер разделяет интерпретацию описания, построение формального документа и его проверку. На рисунке \ref{fig:method-pipeline} показаны основные этапы.

\begin{figure}[H]
    \centering
    \begin{tikzpicture}[
        node distance=0.25cm,
        stage/.style={draw,rounded corners=3pt,minimum width=2.35cm,minimum height=0.95cm,align=center,font=\scriptsize},
        arrow/.style={-{Stealth[scale=0.8]},thick}
    ]
        \node[stage,fill=blue!8] (text) {Текстовое\\описание};
        \node[stage,fill=green!8,right=of text] (brief) {ModelBrief\\и scope check};
        \node[stage,fill=green!8,right=of brief] (plan) {Typed\\ComponentPlan};
        \node[stage,fill=yellow!12,right=of plan] (ir) {Canonical\\ModelDocument};
        \node[stage,fill=orange!12,below=of ir] (validate) {Units, topology,\\graphs, smoke};
        \node[stage,fill=red!8,left=of validate] (repair) {Bounded repair};
        \node[stage,fill=blue!8,right=of validate] (result) {Симуляция\\и экспорт};

        \draw[arrow] (text)--(brief);
        \draw[arrow] (brief)--(plan);
        \draw[arrow] (plan)--(ir);
        \draw[arrow] (ir)--(validate);
        \draw[arrow] (validate)--(result);
        \draw[arrow] (validate)--(repair);
        \draw[arrow] (repair.west) -- ++(-0.35,0) |- (ir.west);
    \end{tikzpicture}
    \caption{Схема ограниченной генерации}
    \label{fig:method-pipeline}
\end{figure}

До обращения к LLM проверяется, не требует ли описание заведомо неподдерживаемых возможностей. Пространственные PDE и стохастические уравнения с шумом Винера отклоняются с `UNSUPPORTED\_CAPABILITY`, а не заменяются детерминированной ODE. Явный интервал `final\_time` $\leq$ `initial\_time` также отклоняется до генерации. Для остальных описаний формируется ModelBrief: цель, границы, временной масштаб, уровни, потоки, внешние воздействия и ожидаемые петли.

Если описание не относится к проверенному archetype, следующим этапом формируется ComponentPlan. При непустом валидном плане переменные создаются построителем, а не свободной генерацией JSON. При пустом плане сохраняется путь variables $\rightarrow$ equations, но он подчиняется тем же проверкам.

\section{Каноническое представление}

ModelDocument содержит версию схемы, описание, временные настройки, переменные, lookup-таблицы, размерности, assumptions, provenance, capabilities и производные CLD/SFD. Переменная относится к одному из четырёх типов:

\begin{itemize}
    \item stock --- начальное значение, единицы и, при необходимости, явная производная;
    \item flow --- выражение и направленные source/target;
    \item auxiliary --- вычисляемое выражение без отдельного состояния;
    \item constant --- конечное число или допустимое константное выражение.
\end{itemize}

Выражения хранятся как ограниченный AST, а не как исходный Python-код. Пример записи скорости роста:

\begin{lstlisting}[caption={Фрагмент канонического AST}]
{"op":"mul","args":[
  {"op":"ref","name":"population"},
  {"op":"ref","name":"growth_rate"}
]}
\end{lstlisting}

Допускаются литералы, ссылки, время, арифметические и логические операции, условия, lookup, индексация и белый список функций. Не разрешаются импорты, обращения к атрибутам, lambda, comprehensions, динамическое индексирование и произвольные вызовы. Ограничения узлов и глубины AST уменьшают риск чрезмерно сложных выражений. Парсер DSL нормализует формулу без `eval`.

Производная уровня вычисляется из топологии по формуле (\ref{eq:flow-balance}). Если явная производная и потоки заданы одновременно, их нормализованные формы должны совпадать; иначе модель отклоняется с `DERIVATIVE\_TOPOLOGY\_MISMATCH`. Нулевые начальные значения сохраняются, а корректное нулевое равновесие не заменяется искусственным притоком.

\section{Типизированные компоненты}

ComponentPlan задаёт ограниченный набор механизмов: `first\_order`, `conditional\_inflow`, `queue`, `scheduled\_flow`, `delay` и `polynomial\_ode`. Каждый тип имеет обязательные поля и ограничения. Например, очередь содержит начальный запас, период поступления, время обслуживания и максимальную скорость обработки. Delay хранит исходный и начальный delayed rate, чтобы не создавать неопределённое состояние.

Для коэффициента перехода от исходной величины $x$ к целевой $y$ единицы вычисляются как

\begin{equation}
    [k]=[y][x]^{-1}[t]^{-1}.
    \label{eq:coefficient-units}
\end{equation}

Композиция компонентов не сводится к объединению списков. Дублирование расхода задержки с той же временной константой отклоняется, поскольку оно способно дважды вычесть один поток. Полностью совпадающие компоненты дедуплицируются по каноническому JSON, но конфликтующие идентификаторы не объединяются.

Повторяющиеся уровни объединяются только при совпадении начальных значений, единиц, производных и ограничений. Исправление duplicate flows выполняется только при доказанной идемпотентности соответствующего адаптера.

\section{Сохранение явных ОДУ}

Если в описании задано $\frac{dX}{dt}=f(X,Y,t)$, свободная декомпозиция LLM может изменить коэффициенты. Поэтому RHS разбирается безопасным parser только на верхнеуровневые слагаемые. Внешний числовой или именованный коэффициент выносится, а внутренние скобки сохраняются.

Например,

\begin{equation}
    \frac{dx}{dt}=a\left(x-\frac{x^3}{3}-y+b\right)-cy
\end{equation}

не раскрывается в сумму без контроля. Каждое слагаемое сохраняет свой знак, а использованные параметры остаются в плане только при реальном присутствии в выражении.

Для оценки эквивалентности используется численная проверка на четырёх наборах значений. Пусть $E_s$ --- исходное выражение, а $E_g$ --- построенная производная. На наборе присваиваний $\mathcal U$ проверяется

\begin{equation}
    \left|E_s(u)-E_g(u)\right|\leq 10^{-8}+10^{-7}|E_s(u)|,\qquad u\in\mathcal U.
    \label{eq:numerical-equivalence}
\end{equation}

Это не символьное доказательство эквивалентности. При неподдерживаемом или неоднозначном выражении результат помечается unknown, а не признаётся успешным. Явно заданные общие units состояний имеют приоритет над догадкой LLM о физических единицах вроде `length`. Аналогично, строковое `"1/time"` нормализуется существующим strict units parser в `{"time": -1}` без принятия malformed типов.

\section{Проверка и ограниченное исправление}

Проверка включает:

\begin{itemize}
    \item существование и уникальность идентификаторов;
    \item конечность литералов и допустимость начальных выражений;
    \item размерность арифметики, функций, условий и потоков;
    \item существование хотя бы одного уровня там, где есть потоки;
    \item направление source/target и соответствие производной потокам;
    \item отсутствие запрещённых auxiliary-циклов;
    \item соответствие CLD/SFD зависимостям выражений с помощью NetworkX \cite{networkx2008};
    \item конечность и полный горизонт контрольной симуляции.
\end{itemize}

При ошибке формируется список диагностик с кодом, путём и сообщением. Repair получает этот список и канонический документ. Для units-only ошибок типы и топология сохраняются; topology-изменения разрешены только при соответствующих классах диагностики. Неудачный ответ repair расходует ограниченный бюджет, но не заменяет предыдущую корректную часть документа на `null`.

\begin{algorithm}[H]
    \caption{Ограниченная проверка и исправление}
    \label{alg:bounded-repair}
    \begin{algorithmic}[1]
        \Require План $P$, бюджет $B>0$
        \State $M\gets$ детерминированно построить или нормализовать $P$
        \State $E\gets$ Validate($M$)
        \While{$E\neq\varnothing$ \textbf{and} $B>0$}
            \State $R\gets$ Repair($M,E$) с сохранением ограничений
            \If{$R$ не проходит схему}
                \State $B\gets B-1$
                \State продолжить с исходным $M$
            \Else
                \State $M\gets R$
                \State $E\gets$ Validate($M$)
                \State $B\gets B-1$
            \EndIf
        \EndWhile
        \If{$E\neq\varnothing$}
            \State \Return ошибка с диагностикой и числом попыток
        \Else
            \State выполнить smoke-симуляцию и экспорт
            \State \Return $M$
        \EndIf
    \end{algorithmic}
\end{algorithm}

Отказ после исчерпания бюджета считается предпочтительнее возврата неправильной модели. При этом отказ сам по себе не доказывает качество системы: он должен анализироваться по причинам и частоте.

\section{Симуляция и экспорт}

Основной движок использует PySD \cite{pysd2021}. Python-файл создаётся доверенным кодогенератором из ModelDocument, а не LLM. Поддерживаются Euler, RK4, RK45, DOP853, LSODA и BDF в соответствующих подмножествах. SciPy fallback ограничен доказанно совместимыми скалярными ODE и не применяется к stateful delay или подскриптам, где его семантика могла бы измениться.

Временная сетка проверяется на полный горизонт и конечность всех рядов. Для imported subscripted variables результаты DataArray разворачиваются по элементам dimensions. Поддерживаются экспорт JSON, Python, Vensim, XMILE и CSV. Экспорт не является независимым источником уравнений: все форматы порождаются из того же документа.

Vensim importer ограничен безопасным подмножеством: скалярные уравнения, lookups, задержки, события, dimensions, числовые многомерные и частичные массивы, element-specific stock derivatives, subranges, aliases и `:EXCEPT:`. Поддерживаются expression/stock/local/nested macros с ограничением глубины и размера. Рекурсия, циклические локальные определения и неверное число аргументов отклоняются. DTD, внешние сущности и произвольные include не принимаются \cite{vensim2024,xmile2015}.

\section{Воспроизводимость}

Каждый benchmark-case может сохранять prompt, ответы LLM, ModelBrief, ComponentPlan, документ, smoke, графы, ошибки и итоговые метрики. Отчёт содержит corpus hash, pipeline hash, модель, repair budget, concurrency, среднее время и targets. Повторный запуск одного case очищает его прежние JSON-стадии, не затрагивая другие cases.

Для метрик используются следующие определения:

\begin{itemize}
    \item \textbf{executable} --- документ принят, симуляция запущена и завершена;
    \item \textbf{semantic} --- ожидаемый archetype совпадает, rule-based semantic issues отсутствуют, units/topology/graph проверки пройдены;
    \item \textbf{behavior} --- выполнен semantic, все значения конечны, горизонт полный, динамика соответствует ожиданию, а специальный behavior oracle пройден.
\end{itemize}

Semantic здесь является алгоритмической проверкой, а не автоматическим доказательством содержательной адекватности. Для экспертной оценки требуется отдельный протокол аннотации. Повторные генерации одного описания служат для оценки вариативности и не должны рассматриваться как независимые описания в статистическом анализе.
